\clearpage
\fchapter{Objetivos}

El objetivo de esta práctica es planificar y realizar el montaje del control de temperatura de la maqueta de horno Alecop MT-542 mediante un controlador de temperatura Omron E5EC, de forma que el horno se mantenga en un punto de consigna y señalice su estado mediante alarmas. Para ello se plantean los siguientes objetivos concretos:

\begin{itemize}
	\item Planificar la instalación: identificar los bornes del controlador a partir de su documentación técnica y elaborar el esquema eléctrico del conjunto.
	\item Realizar el conexionado del controlador: alimentación, sonda Pt100 a tres hilos, salida de control hacia la resistencia calefactora de la maqueta y salidas de alarma.
	\item Configurar el controlador desde su panel frontal: tipo de entrada, control PID y ajuste automático de sus parámetros mediante autotuning.
	\item Configurar cuatro alarmas de temperatura, de valor absoluto y de desviación respecto al punto de consigna, y utilizar sus salidas de relé para gobernar los tres pilotos de un módulo de señalización y el ventilador de la maqueta.
	\item Comprobar el correcto funcionamiento del conjunto, verificando que la temperatura se estabiliza en la consigna y que cada alarma actúa en el rango previsto.
\end{itemize}
